% year: תשפ"ג
% semester: ב
% moed: א
% note: ד"ר פיטר סמובול
% source: exams/מבחנים/תשפג/מועד א תשפג סמסטר ב.pdf
% number: 2א
% categories: antiderivatives

\begin{question}
(שאלה חובה)
פתרו את האינטגרל הבא:
$\displaystyle\int\frac{(x^2+1)\cdot dx}{x^3-x}$
\end{question}

\begin{hint}
פרקו לשברים חלקיים: $x^3-x=x(x-1)(x+1)$.
\end{hint}

\begin{solution}
$x^3-x=x(x-1)(x+1)$, ודרגת המונה קטנה מדרגת המכנה. נפרק:
\[ \frac{x^2+1}{x(x-1)(x+1)}=\frac{A}{x}+\frac{B}{x-1}+\frac{C}{x+1}. \]
בשיטת הכיסוי: $A=\frac{0+1}{(0-1)(0+1)}=-1$, $B=\frac{1+1}{1\cdot 2}=1$, $C=\frac{1+1}{(-1)(-2)}=1$. לכן
\[ \int\frac{x^2+1}{x^3-x}dx=-\ln|x|+\ln|x-1|+\ln|x+1|+C=\boxed{\ln\left|\frac{x^2-1}{x}\right|+C}. \]
\end{solution}
